% SPDX-License-Identifier: LPPL-1.3c
% Copyright (C) 2026 Christophe Jorssen
% Author and Current Maintainer: Christophe Jorssen <christophe.jorssen@gmail.com>
% LPPL maintenance status: maintained
% This file is part of luafemm. See LICENSE and LICENSES.md for its terms.
\section{Effective windings and nonlinear fundamental fields}
\label{sec:harmonic-material-models}

A fine-wire winding and a saturated core ask two different questions. In the
winding, current crowds inside each conductor and neighbouring turns induce
additional currents. In the core, the slope of the B--H curve changes during a
cycle. The models in this chapter address these effects separately. Both use
ordinary material styles, the existing circuit keys and the existing field,
profile and measurement commands. Their applicability limits are part of their
interfaces: selecting an unsupported combination raises an error.

\subsection{Replacing a thick wire by Litz wire}
Suppose a winding uses twenty turns of round copper wire. Its copper diameter
is 0.4 mm. At low frequency, a uniform current-density model is sufficient;
at higher frequency, it misses skin and proximity effects. Describe the wire
in the material style and connect the two signed winding sections to a circuit:
\begin{codeexample}[code only]
\tikzset{femm/materials/wire/.style={mu r=1,conductivity=58000000,
  wire model=round,wire diameter=.4}}
\end{codeexample}
The diameter is the conducting diameter, excluding insulation. Do not enter a
separate |winding fill|: the copper area follows from turns and diameter, and
the available winding area follows from the captured geometry. This gives the
same physical filling factor when the mesh is refined.

Now replace each wire by sixteen insulated 0.1 mm strands, perfectly transposed
so that each strand sees the same average environment. The copper area per turn
is unchanged. Change the style to |wire model=litz,wire diameter=.1,strands=16|.
The following complete comparison starts with a short frequency list and a
moderate mesh. Refine the winding and surrounding air when the diagram and
circuit have settled; meshing the wire diameter itself is unnecessary.

\luafemmexamplebegin{wire-windings}

% manual example: wire-windings/1

\begin{codeexample}[setup code]
\newbox\Calculation
\tikzset{
  femm/materials/wire/.style={mu r=1,conductivity=58000000,
    wire model=\WireModel,wire diameter=\Diameter,strands=\Strands},
  experiment/.style={femm/problem={frequency=\Frequency,depth=20,
    xmin=-8,xmax=8,ymin=-6,ymax=6,mesh size=.5},
    femm/circuit={name=drive,current=2}},
  /femm/region/coil/.style={material=wire,circuit=drive,mesh size=.25}}
\newcommand\Geometry{
  \path[femm/region={coil,turns=-20}] (-4.7,-2) rectangle (-3.3,2);
  \path[femm/region={coil,turns=20}] (3.3,-2) rectangle (4.7,2);}
\foreach \WireModel/\Diameter/\Strands in {round/.4/1,litz/.1/16}{
  \foreach \Frequency in {100,1000,5000,10000,50000,100000}{
    \setbox\Calculation=\hbox{
      \begin{tikzpicture}[experiment,
        femm/cache={file=\jobname-\WireModel-\Frequency,mode=auto}]
        \Geometry
        \femmmeasure[name=\WireModel-\Frequency,parameter=\Frequency]
      \end{tikzpicture}}}}
\def\Frequency{100000}
\end{codeexample}

The following picture uses the preceding definitions.

% manual example: wire-windings/2

\begin{codeexample}[width=0pt]
\foreach \WireModel/\Diameter/\Strands/\Title in
  {round/.4/1/{Round wire},litz/.1/16/{Litz wire}}{
  \begin{tikzpicture}[experiment,scale=3.2,
    femm/cache={file=\jobname-\WireModel-\Frequency,mode=auto}]
    \Geometry
    \filldraw[fill=orange!20,draw=orange!65!black] (-4.7,-2) rectangle (-3.3,2);
    \filldraw[fill=orange!20,draw=orange!65!black] (3.3,-2) rectangle (4.7,2);
    \pic[femm/lines=15,femm/phase=0,femm/arrow spacing=12mm] {femm field};
    \node[above] at (0,6) {\Title, $100\,\mathrm{kHz}$};
    \node at (-4,0) {$\otimes$}; \node at (4,0) {$\odot$};
    \node[below,font=\small] at (0,-6)
      {$P=\femmvalue[component=total_loss]{\WireModel-100000}$};
  \end{tikzpicture}}
\end{codeexample}

% manual example: wire-windings/3

\begin{codeexample}[setup code]
\def\Round{round-100,round-1000,round-5000,round-10000,round-50000,round-100000}
\def\Litz{litz-100,litz-1000,litz-5000,litz-10000,litz-50000,litz-100000}
\end{codeexample}

The following picture uses the preceding definitions.

% manual example: wire-windings/4

\begin{codeexample}[width=5.2cm]
\begin{tikzpicture}
\begin{semilogxaxis}[width=7.5cm,height=6cm,grid=major,
  xlabel={Frequency (Hz)},ylabel={Total winding loss (W)},
  legend style={font=\small,at={(.03,.97)},anchor=north west}]
  \femmresultplot[results={\Round},component=total_loss,
    plot options={blue!70!black,thick,mark=*}]
  \addlegendentry{Round, $d=0.4$ mm}
  \femmresultplot[results={\Litz},component=total_loss,
    plot options={red!70!black,dashed,thick,mark=square*}]
  \addlegendentry{Litz, $16\times0.1$ mm}
\end{semilogxaxis}
\end{tikzpicture}
\end{codeexample}

% manual example: wire-windings/5

\begin{codeexample}[width=5.2cm]
\begin{tikzpicture}
\begin{loglogaxis}[width=7.5cm,height=6cm,grid=major,
  xlabel={Frequency (Hz)},ylabel={Round-wire loss terms (W)},
  legend style={font=\small,at={(.03,.97)},anchor=north west}]
  \femmresultplot[results={\Round},component=wire_transport_loss,
    plot options={blue!70!black,thick,mark=*}]
  \addlegendentry{Transport term}
  \femmresultplot[results={\Round},component=proximity_loss,
    plot options={red!70!black,dashed,thick,mark=square*}]
  \addlegendentry{Proximity term}
\end{loglogaxis}
\end{tikzpicture}
\end{codeexample}

\luafemmexampleend{wire-windings}

The total loss includes both terms shown in the second plot. The transport term
includes an effective-medium correction and need not, on its own, exceed the DC
loss. Compare the total losses of geometries with the same copper area and
excitation. The model excludes end turns, capacitive effects and imperfect
transposition. Adding more strands does not automatically make a physically
manufacturable Litz bundle or establish the model's accuracy for that bundle.

\subsection{Wire material keys}
\label{sec:wire-keys}
These keys belong to |/tikz/femm/materials|. Define them in material styles
before capturing any affected region. Empty numeric values inherit catalogue
properties, then the native defaults. A material style is evaluated separately
for each region; a material can therefore be shared by winding sections with
different areas and derived fills. Ordinary TeX grouping scopes style changes.

\begin{key}{/tikz/femm/materials/wire model=\meta{choice} (initially inherit)}
A choice handler accepting |inherit|, |none|, |round|, |stranded|, |litz|, |square|, |cca10| and |cca15|.
|inherit| selects supported catalogue wire data at positive frequency,
otherwise |none|. Static catalogue selection uses bulk conductivity and unit
permeability; an explicitly selected nontrivial wire law still requires AC. With |none|,
the existing uniform stranded conductor or resolved solid-conductor model is
unchanged. |round| models one round wire per turn. |stranded| treats an
untransposed bundle as one equal-copper-area round wire. |litz| uses the
individual strand radius and assumes ideal transposition. These are
homogenisation choices, not geometric generators.

A nontrivial wire model requires positive |frequency|, conductivity greater
than zero, $\mu_x=\mu_y=1$, no B--H curve, no coercivity, no sheet model and no
intrinsic loss angle. Every wire region must belong to a named circuit and have
|conductor=auto| or |stranded|. Both current and voltage drive, and series and
parallel branches, use the coupled circuit solver. Positive |depth| is required
in planar models; axisymmetric integrals use the actual ring circumference.
Captured operating points and the axial diffusion formulation are excluded.

Leave |winding fill=1| and |conductor length=0|: these values request geometry-
derived quantities, not physically full copper occupancy. Solid-current
constraints and independent source-density regions are rejected. Static use of
a selected wire model is rejected; use an ordinary conductive material with an
explicit winding fill for a DC model. The additional square and copper-clad aluminium models are described below.
\end{key}

\begin{key}{/tikz/femm/materials/wire diameter=\meta{millimetres} (initially empty)}
Positive finite conducting diameter (square side for |square|) in millimetres, independent of |unit|. For
|stranded| and |litz| this is the diameter of one strand. There is no guessed
positive default; a nontrivial wire model requires this value or supported
catalogue data. The native value zero is valid only with |wire model=none|.
PGF expressions are accepted, for example |wire diameter=.4/sqrt(16)|.
\end{key}

\begin{key}{/tikz/femm/materials/strands=\meta{integer} (initially empty)}
Number of strands per electrical turn. Native default one; accepts positive
integers, including integer-valued PGF expressions. |round| and |square| require one.
The circuit's signed |turns| still counts turns of the complete bundle, not
individual strands. |wire model=none| requires the default strand count and
zero wire diameter, preventing silently ignored wire data.
\end{key}

For winding area $S$, signed turn count $N$, strand count $n$ and strand radius
$a=d/2$ in metres, the copper fraction is
\[
  \eta=\frac{|N|n\pi a^2}{S},\qquad
  0<\eta\leq\frac{\pi}{\sqrt{12}}.
\]
The upper bound is the round-wire hexagonal packing limit; insulation normally
requires a smaller fraction. Overfilled regions fail after their mesh areas
are available. The effective skin radius is $a$ for round/Litz wire and
$a\sqrt n$ for untransposed strands. Thus the round and untransposed equal-area
models coincide for the example's parameters, while the Litz model differs.

Catalogue |LamType=3| round-wire records activate this law in harmonic solves;
|LamType=4| and |5| conversions use the corresponding strand models when present.
Types 6, 7 and 8 select square, CCA10 and CCA15 respectively.
The wire dimensions remain catalogue data. Unsupported records fail explicitly.

\subsection{What is integrated, and what is drawn?}
\label{sec:wire-losses}
The approximation supplies a complex relative reluctivity $\nu_r^*$ and an
internal effective resistivity $\rho^*$. With peak phasors,
\[
 p_{\rm transport}=\tfrac12\Re(\rho^*)|\widehat J|^2,
 \qquad
 p_{\rm proximity}=\frac{\omega}{2\mu_0}\Im(\nu_r^*)|\widehat{\mathbf B}|^2.
\]
Here $\widehat J=NI/S$ is the \emph{macroscopic} section current density, not the
current density inside the copper. Its transport voltage contribution is
$Z_{\rm local}I$, with $Z_{\rm local}=N^2\ell\rho^*/S$ for section length $\ell$.
The circuit also includes $\mathrm i\omega\lambda$ from the computed potential.
The effective resistivity includes the local-energy correction used in FEMM's
wire formulation, so this energy is not counted again in both terms.

\begin{center}
\begin{tabular}{lp{.64\linewidth}}
\toprule Measurement component & Meaning in a harmonic integral\\
\midrule
|wire_transport_loss| & Integral of $p_{\rm transport}$, in W.\\
|proximity_loss| & Integral of $p_{\rm proximity}$, in W.\\
|wire_energy| & Signed local correction $\int\Im(\rho^*)|\widehat J|^2/(4\omega)\,dV$, in J.\\
|joule_loss| & Ordinary solid/stranded losses plus both wire loss terms, exactly once.\\
|eddy_loss| & Resolved bulk eddy-current contribution only; zero inside a homogenised wire region.\\
|magnetic_loss| & Intrinsic and sheet losses, as in Section~\ref{sec:losses}; excludes wire losses.\\
|total_loss| & |joule_loss| plus |magnetic_loss|.\\
\bottomrule
\end{tabular}
\end{center}
These quantities use |\femmmeasure|, |\femmvalue| and |\femmresultplot| with the
usual selection, formatting and plot options. All three new wire components
are zero in ordinary planar/axisymmetric harmonic materials. |energy| includes
|wire_energy|; for these linear models |coenergy=energy|. The terminal
|resistance| remains the geometrically derived DC resistance, while
|impedance_real| and |impedance_imaginary| include frequency-dependent transport
and field contributions. A port's power is not a regional loss if other ports
exchange energy with it; compare the sum of terminal powers with total losses.

Field lines and B/H profiles show the homogenised magnetic field. Current
profiles and glyphs show the macroscopic transport current; |Jeddy=0| inside
a wire region. They cannot display microscopic skin/proximity loops without
explicitly resolving the individual wires. Lorentz forces in the homogenised
wire material are rejected; a surrounding-air Maxwell selection remains subject
to the force chapter's usual geometry and material checks.

The effective-medium construction follows David Meeker's
\href{https://www.femm.info/doku/doku.php?id=proximityloss}{skin/proximity model}.
The fitted coefficients and local-energy convention were audited against
FEMM/xfemm's |GetFillFactor| implementation. The Lua runtime evaluates the
mathematical model without a native dependency. Agreement with xfemm checks
implementation of the same approximation; it does not validate arbitrary real
winding constructions. See the validation report for independent meshes and
measured discrepancies.

\subsection{A saturated core: choose the approximation explicitly}
\label{sec:fundamental-law}
A static B--H curve is not, by itself, a complete AC material description.
If hysteresis and eddy currents are negligible and only the fundamental field
is required, an explicitly selected single-harmonic approximation is useful.
Use a single-valued isotropic table and |harmonic law=fundamental|. All
excitations must be real, phase zero, prescribed currents or real boundary data.
This assumption fixes the direction of the local oscillating flux; it does not
represent a rotating elliptical flux locus.

\begin{key}{/tikz/femm/materials/harmonic law=\meta{choice} (initially inherit)}
Choice handler accepting |inherit|, |standard|, |fundamental| and |equivalent|. |inherit|
uses the material's setting, otherwise |standard|. The standard law retains
linear harmonic behaviour or the separately requested operating-point response.
It never silently converts a B--H curve into a large-signal harmonic law.

Select |fundamental| for a positive-frequency model with an isotropic |bh| table.
The |interpolation| problem option must be |linear| (the default).
This law projects $H(B)$ assuming sinusoidal
$B$, then solves the real nonlinear fundamental balance. The DC B--H data are
preserved. No loss angle, lamination, permanent-magnet coercivity, operating
point, wire homogenisation, voltage drive, parallel circuit, isolated solid or
resolved eddy current may be combined with this approximation. Conductive coil
materials are allowed only with a stranded interpretation; their ordinary
transport loss remains available. All circuit and current-density phases must
be exactly zero. The planar and axisymmetric vector-potential formulations
are supported; the axial out-of-plane formulation is excluded.

The standard Newton tolerance and iteration limit govern convergence. Failure
does not produce a saved solution. Force/torque measurements are rejected for
this approximation because neglected harmonics can contribute to mechanical
loads. Use the documented static or linear harmonic force models instead.
\end{key}

\luafemmexamplebegin{nonlinear-fundamental}

% manual example: nonlinear-fundamental/1

\begin{codeexample}[setup code]
\newbox\Calculation
\tikzset{
  femm/materials/core/.style={bh={0/0,1/100,1.5/500,2/5000},
    harmonic law=\Law},
  experiment/.style={femm/problem={frequency=\Frequency,depth=20,
    xmin=-8,xmax=8,ymin=-6,ymax=6,mesh size=.5},
    femm/circuit={name=drive,current=\Current}},
  /femm/region/coil/.style={material=copper,circuit=drive,winding fill=.6}}
\newcommand\Geometry{
  \path[even odd rule,femm/region={material=core,mesh size=.25}]
    (-3,-4) rectangle (3,4) (-1.5,-2.5) rectangle (1.5,2.5);
  \path[femm/region={coil,turns=-20}] (.8,-2) rectangle (1.3,2);
  \path[femm/region={coil,turns=20}] (3.7,-2) rectangle (4.2,2);}
\foreach \Case/\Law/\Frequency in {ac/fundamental/50,dc/standard/0}{
  \foreach \Index/\Current in {a/.1,b/.3,c/1,d/3,e/10,f/30}{
    \setbox\Calculation=\hbox{
      \begin{tikzpicture}[experiment,
        femm/cache={file=\jobname-\Case-\Current,mode=auto}]
        \Geometry
        \femmmeasure[name=\Case-\Index,quantity=circuit,circuit=drive,parameter=\Current]
        \path[femm/profile={name=profile-\Case-\Index,samples=61}]
          (1.6,0)--(2.9,0);
      \end{tikzpicture}}}}
\def\Law{fundamental}\def\Frequency{50}\def\Current{10}
\end{codeexample}

The following picture uses the preceding definitions.

% manual example: nonlinear-fundamental/2

\begin{codeexample}[width=0pt]
\begin{tikzpicture}[experiment,scale=6,
  femm/cache={file=\jobname-ac-10,mode=auto}]
  \Geometry
  \filldraw[even odd rule,fill=blue!10,draw=blue!50!black]
    (-3,-4) rectangle (3,4) (-1.5,-2.5) rectangle (1.5,2.5);
  \filldraw[fill=orange!25,draw=orange!65!black] (.8,-2) rectangle (1.3,2);
  \filldraw[fill=orange!25,draw=orange!65!black] (3.7,-2) rectangle (4.2,2);
  \pic[femm/lines=11,femm/arrow spacing=12mm] {femm field};
  \draw[red!75!black,very thick] (1.6,0)--(2.9,0);
  \node[above] at (0,5) {Fundamental approximation, $10$ A peak};
  \node at (1.05,0) {$\otimes$}; \node at (3.95,0) {$\odot$};
\end{tikzpicture}
\end{codeexample}

% manual example: nonlinear-fundamental/3

\begin{codeexample}[width=0pt]
\begin{tikzpicture}
\begin{semilogxaxis}[width=7.5cm,height=6cm,grid=major,
  xlabel={Peak current (A)},ylabel={Flux linkage (Wb turn)},
  legend style={font=\scriptsize,at={(.03,.97)},anchor=north west}]
  \femmresultplot[results={ac-a,ac-b,ac-c,ac-d,ac-e,ac-f},
    component=linkage_real,plot options={blue!70!black,thick,mark=*}]
  \addlegendentry{Fundamental at 50 Hz}
  \femmresultplot[results={dc-a,dc-b,dc-c,dc-d,dc-e,dc-f},
    component=linkage,plot options={red!70!black,dashed,thick,mark=square*}]
  \addlegendentry{Static at peak current}
\end{semilogxaxis}
\end{tikzpicture}
\end{codeexample}

% manual example: nonlinear-fundamental/4

\begin{codeexample}[width=0pt]
\begin{tikzpicture}
\begin{axis}[width=7.5cm,height=6cm,grid=major,
  no markers,yticklabel style={/pgf/number format/fixed,/pgf/number format/precision=3},
  xlabel={$x$ in the right limb (mm)},ylabel={$B_y$ (T)},
  legend style={font=\scriptsize,at={(.5,1.03)},anchor=south}]
  \femmplot[profile=profile-ac-e,component=By,abscissa=x,representation=real,
    plot options={blue!70!black,thick}]
  \addlegendentry{Fundamental amplitude}
  \femmplot[profile=profile-dc-e,component=By,abscissa=x,
    plot options={red!70!black,dashed,thick}]
  \addlegendentry{Static at peak current}
\end{axis}
\end{tikzpicture}
\end{codeexample}

\luafemmexampleend{nonlinear-fundamental}

The comparison uses the same geometry and peak current in both solves. Static
flux at peak current and the retained harmonic amplitude generally differ once
saturation matters. Neither curve is a reconstructed transient solution. Begin
with a few amplitudes and a coarse mesh; then refine the core, interfaces and
surrounding air and check both flux linkage and local field profiles. A current
sweep changes the physical solution and must not reuse a frozen cache from a
different amplitude.

\subsection{Projection, energy and limits of interpretation}
Write $b=|\widehat{\mathbf B}|$. The projected constitutive law and its tangent are
\[
 H_1(b)=\frac4\pi\int_0^{\pi/2}H(b\cos\theta)\cos\theta\,d\theta,
 \qquad
 H'_1(b)=\frac4\pi\int_0^{\pi/2}H'(b\cos\theta)\cos^2\theta\,d\theta.
\]
Each linear B--H segment is integrated analytically between its cosine crossing
angles. The continuation beyond the last point has vacuum slope, matching the
static piecewise-linear law. Avoid relying on that continuation to replace
measured saturation data. The consistent tangent drives the damped Newton
solver; this avoids fitting a new table at every trial field.

The reported cycle-average energies are
\[
 \overline w=\frac2\pi\int_0^{\pi/2}
       \left[\int_0^{b\cos\theta}H(s)\,ds\right]d\theta,
 \qquad
 \overline w'=\tfrac12 bH_1(b)-\overline w.
\]
In general they are unequal, and neither can be replaced by $bH_1/4$.
The |energy| and |coenergy| measurements use these expressions. In the lossless
current-driven limit, reactive terminal power satisfies
$Q=\omega\int(\overline w+\overline w')\,dV$.
H/B profiles and phase snapshots use the retained fundamental phasors. In
particular, an H snapshot displays $H_1\cos\omega t$, not a reevaluation of
the full instantaneous law $H(b\cos\omega t)$.
Hysteresis loss is zero because the underlying law is single-valued and
reversible. A nonlinear waveform can contain higher harmonics even when its
fundamental is known; this solver does not compute them.

This projection assumes sinusoidal B and projects H(B). FEMM's effective
large-signal curve construction instead projects B(H) under its own assumptions;
the two are not generally equivalent. Accordingly, nonlinear validation uses
independent integration of the original B--H curve, manufactured uniform fields,
linear limits, nonlinear residuals and energy identities, rather than claiming
numerical equivalence to FEMM's large-signal solver.

Native cache revision 25 fingerprints wire dimensions, strand interpretation
and harmonic-law choice. Derived fill factors, complex coefficients and network
states are rebuilt from saved geometry and potentials. Changing these physical
inputs invalidates the solution; only compatible geometry may reuse its mesh.
Supported linear, sheet and wire harmonic FEM/ANS interchange is described in
Section~\ref{sec:interchange}. Fundamental-law export remains unsupported.

\subsection{Square and copper-clad aluminium conductors}
The |square| choice interprets |wire diameter| as the side $d$ of one insulated
square conductor. Its area is $d^2$, |strands=1| is required and the geometrical
filling fraction may not exceed one. An equivalent foil stack, with linear fill
$\sqrt\eta$, supplies both the field reluctivity and the terminal resistivity.
This is a homogenized array, not a resolved square-wire current distribution.
Use a resolved conducting region when corner crowding itself is of interest.

The |cca10| and |cca15| choices describe round copper-clad aluminium conductors
with respectively 10\% and 15\% copper cross-sectional fractions. Conductivity is
the \emph{bulk DC conductivity of the composite}, not the conductivity of copper.
The fitted skin/proximity laws require $0.05\leq\eta\leq\pi/\sqrt{12}$.
The stated diameter excludes insulation; the strand count counts equal wires.
Their fit coefficients follow the FEMM 5 conductor models. The square model uses
one consistent diffusion frequency for field and terminal calculations; FEMM 5's
preprocessor and postprocessor use different square-wire factors. Consequently,
exact numerical identity with both native paths is not claimed.

For example, use the existing winding declaration with the following style.
The conductor area and turn count must fit inside the captured winding region.
\begin{codeexample}[code only]
\tikzset{femm/materials/composite wire/.style={mu r=1,
 conductivity=38000000,wire model=cca10,wire diameter=.4,strands=1}}
\end{codeexample}
All winding transport/proximity components and terminal power balances retain
the definitions above. These models require no native library at runtime.

\subsection{An equivalent nonlinear harmonic law with losses}
\label{sec:equivalent-harmonic}
Choose |harmonic law=equivalent| for an isotropic, piecewise-linear B--H curve
with a common |loss angle| in both directions. For a peak vector phasor define
\[
 B_{\rm eq}=\sqrt{|\widehat B_x|^2+|\widehat B_y|^2},\qquad
 H_1(b)=\frac4\pi\int_0^{\pi/2}H(b\cos\theta)\cos\theta\,d\theta,
 \qquad \nu_r^*=\mu_0\frac{H_1(B_{\rm eq})}{B_{\rm eq}}e^{i\delta}.
\]
The zero-amplitude limit uses the initial slope. This is the fundamental of a
linearly polarized sinusoid followed by a prescribed phase lag. For an ellipse,
$B_{\rm eq}$ is $\sqrt2$ times vector RMS, not the peak vector magnitude.
No harmonics, remanence history, minor loops or irreversible demagnetization are
reconstructed. The loss-angle convention is passive for the supported angles.

This law accepts fixed planar or axisymmetric fields and prescribed complex
stranded current densities. It excludes conductivity, wire/sheet homogenization,
permanent magnets, coupled voltage circuits, networks, flux controllers,
operating points, motion and time integration. Declare ideal impressed-current
regions with a nonconducting material when using this approximation. Do not mix
|fundamental| and |equivalent| materials in the same problem.

The existing problem |tolerance| and |max iterations| keys also control a damped
coefficient iteration. Both the inner linear solve and the constitutive
fixed-point defect must converge. Strong saturation may require a larger
iteration budget; failure is reported instead of drawing an unconverged field.
Harmonic losses use the final complex reluctivity. Reported harmonic energy is
the equivalent quadratic quantity $\frac14\operatorname{Re}(\widehat{\mathbf B}
\cdot\overline{\widehat{\mathbf H}})$, not the exact cycle-averaged nonlinear stored
energy. Force integrals are unavailable with this effective constitutive law.

The following deliberately simple imposed-flux patch makes the units and
representation visible before introducing a complicated core geometry.
\begin{codeexample}[width=0pt]
% Generic body: a passive equivalent-amplitude nonlinear harmonic core.
\begingroup
\tikzset{femm/materials/lossy core/.style={
 bh={0/0,1/100,1.5/500,2/5000},harmonic law=equivalent,loss angle=15}}
\tikzpicture[femm/problem={frequency=50,depth=10,
 xmin=0,xmax=2,ymin=0,ymax=2,mesh size=.25,max iterations=150},
 femm/boundary={side=all,type=dirichlet,y coefficient=1.8},scale=15]
 \path[femm/region={material=lossy core}] (0,0) rectangle (2,2);
 \pic[femm/visual={component=H,representation=amplitude,name=core-H}] {femm scalar map};
 \pic[femm/lines=5] {femm field};
 \draw (0,0) rectangle (2,2);
\endtikzpicture
\tikzpicture\femmcolorbar[name=core-H,height=3cm]\endtikzpicture
\par
\endgroup
\end{codeexample}
